\chapter{总结与展望}
\section{全文总结}
自动驾驶技术作为人工智能和机器人领域的一个重要研究方向，近年来已经取得了显著的进展。自动驾驶赛车是一个具有挑战性的测试平台，能够在高速、动态的环境中验证自动驾驶算法的性能。高速运动状态下赛车的动力学模型复杂度高，难以估计，给控制器的精准预测和控制带来了挑战；环境复杂且变化快，要求算法具备强大、高效的感知控制能力；同时，赛车的安全性要求高，如何平衡安全性与高速运动性能也是一个研究难点。针对这些挑战，本文主要关注赛车的安全控制、端到端竞速以及端到端高速导航三个研究问题，所应用的场景包括已知动态环境的安全竞速、未知封闭赛道的视觉竞速以及未知复杂环境的点到点导航。研究成果可以扩展到极限工况下的自动驾驶系统设计，提升自动驾驶系统在复杂环境中的安全性和性能。本文的主要贡献包括：

\begin{enumerate}
    \item 针对已知动态环境中的安全竞速问题，本文首先建立了赛车的动力学模型，并用魔术公式来描述车胎的力学特性。通过收集赛车在不同工况下的运动数据，使用系统辨识方法对模型参数进行估计和优化。该模型能够准确地描述赛车的运动行为，并为后续的控制算法设计提供了基础。基于该模型，本文首先介绍了控制障碍函数方法，该方法将车辆约束在安全区域内，确保赛车在高速运动过程中避免碰撞。然后，本文提出了一种基于模型预测控制的安全竞速方法，将控制障碍函数作为约束求解优化问题，从而实现安全的高速竞速。最后，本文在仿真、实车环境中对所提出的方法进行了验证，结果表明该方法能有效地提高赛车的安全性，同时保持较高的竞速性能。
    \item 针对未知封闭赛道中的视觉竞速问题，本文提出了一种两阶段基于强化学习的端到端竞速方法。第一阶段，使用具有特权信息的强化学习算法训练一个教师策略，该策略能够在已知环境中实现高性能竞速。第二阶段，使用知识蒸馏的方法将教师策略迁移到学生策略，使其能够在未知环境中通过视觉输入实现竞速。该方法提高了策略在高维视觉输入空间的训练效率，解决了部分可观环境中的最优控制问题。为了提高学生策略的泛化能力，本文引入了域随机化技术，在训练过程中对环境参数进行随机化。最后，本文在仿真环境中对所提出的方法进行了验证，与传统控制方法和其他端到端方法相比表现出更优的性能。另外，本文实现了仿真到现实的迁移，并在实车环境中验证了该方法的有效性。
    \item 针对未知复杂环境中的点到点高速导航问题，本文提出了一种基于视觉预测模型的端到端高速导航方法。本文首先设计了一个视觉预测模型并介绍其训练方法，该模型采用了循环状态空间模型的结构，能够从视觉输入中提取环境特征，并预测未来的环境状态。基于该视觉预测模型，本文引入了模型预测路径积分控制方法，通过采样未来的环境状态来优化控制策略。另外针对高速导航问题，本文设计了一个奖励函数，用于评估采样轨迹的质量，从而指导控制策略的优化。该方法解决了强化学习方法在高维视觉输入下的训练效率问题，以及部分可观环境中长距离导航的短视问题。本文在多个仿真环境中进行了验证，并与强化学习等方法进行了比较，在多个指标上表现出明显优势。最后，本文还在实车环境中验证了该方法的有效性，并通过实车数据微调进一步提升了导航性能。
\end{enumerate}
\section{未来研究展望}
本文从安全竞速、视觉竞速和高速导航三个方面研究了自动驾驶赛车的控制问题，并提出了相应的方法和算法。未来的研究可以从以下几个方面进行拓展：
\begin{enumerate}
    \item 在安全竞速方面，未来的研究可以考虑更多动态环境因素的影响，如与其他赛车的博弈行为、赛道路面条件的变化等，以提高算法的鲁棒性和适应性。另外，可以探索更高效的优化算法来求解模型预测控制问题，以满足实时控制的需求。
    \item 在视觉竞速方面，未来的研究可以考虑引入多模态的传感器信息，如激光雷达、毫米波雷达等，以提高环境感知的准确性和鲁棒性。同时，可以探索更先进的强化学习算法，如分层强化学习、多智能体强化学习等，以进一步提升竞速性能。
    \item 在高速导航方面，未来的研究可以考虑引入更多的环境信息，如语义信息、交通规则、人类语言引导等，以提高导航的智能性和交互性。另外，可以探索更高效的视觉预测模型结构，如Transformer等，以提升预测性能和控制效果。
\end{enumerate}

自动驾驶赛车作为一个具有挑战性的测试平台，为自动驾驶技术的研究提供了宝贵的机会。高速赛车控制算法的研究不仅能够提升赛车的性能和安全性，还能够为实际自动驾驶系统的设计提供重要的参考和启示，推动实现更安全、更高效的自动驾驶系统，满足未来智能交通领域的发展需求。